\section*{Bab \ref{ch:g10:heart-lungs-effort} --- Jantung dan Paru-paru Selama Usaha}

\begin{solution}{exo:g10:heart-lungs-effort:1}
$0.6 \times 14 = \qty{8.4}{L/min}$; $2.2 \times 35 = \qty{77}{L/min}$.
\end{solution}

\begin{solution}{exo:g10:heart-lungs-effort:2}
Volume darah yang dipompa satu sisi jantung tiap menit:
$Q = f \times V_s$. $65 \times 0.075 \approx \qty{4.9}{L/min}$.
\end{solution}

\begin{solution}{exo:g10:heart-lungs-effort:3}
Serambi kanan, bilik kanan, arteri paru, lalu paru-paru (peredaran
paru), vena paru, serambi kiri, bilik kiri, aorta, lalu organ
(peredaran sistemik), vena, dan kembali ke serambi kanan.
\end{solution}

\begin{solution}{exo:g10:heart-lungs-effort:4}
Volume tidal \qty{0.5}{L}; kapasitas vital \qty{4.5}{L}; volume sisa
\qty{1.2}{L}.
\end{solution}

\begin{solution}{exo:g10:heart-lungs-effort:5}
Sekitar $220 - 16 = 204$ dan $220 - 60 = 160$ denyut tiap menit.
\end{solution}

\begin{solution}{exo:g10:heart-lungs-effort:6}
\qty{100}{W}: $120 \times 0.105 \approx \qty{12.6}{L/min}$;
\qty{300}{W}: $198 \times 0.115 \approx \qty{22.8}{L/min}$; faktornya
1.8.
\end{solution}

\begin{solution}{exo:g10:heart-lungs-effort:7}
Istirahat: \qty{2.5}{L/min}; olahraga: $0.04 \times 24 \approx
\qty{1.0}{L/min}$. Bagiannya turun dua belas kali lipat, aliran
mutlaknya hanya turun dengan faktor 2.6, karena curahnya sendiri sudah tumbuh.
\end{solution}

\begin{solution}{exo:g10:heart-lungs-effort:8}
$18 \times 0.14 = \qty{2.5}{L/min}$: sekitar tiga perempat dari
\qty{3.45}{L/min} seorang murid terlatih --- dekat ke sana bagi orang
yang lebih kecil.
\end{solution}

\begin{solution}{exo:g10:heart-lungs-effort:9}
Otak tidak dapat menahan bahkan beberapa sekon tanpa oksigen dan
glukosa, sehingga pasokannya dilindungi. Usus dan ginjal dapat
menghentikan pencernaan dan penyaringan selama satu jam tanpa
kerusakan, dan bagiannya dipinjamkan kepada ototnya.
\end{solution}

\begin{solution}{exo:g10:heart-lungs-effort:10}
$V_s = 5000/42 \approx \qty{120}{mL}$, berhadapan dengan sekitar
\qty{70}{mL}: jantungnya yang membesar mengantarkan curah istirahatnya
dengan denyut yang jauh lebih sedikit.
\end{solution}

\begin{solution}{exo:g10:heart-lungs-effort:11}
Perintah otak untuk bergerak, dan antisipasinya terhadap usaha itu,
dikirim ke pusat jantung dan kelenjar adrenalnya sebelum ototnya
bertindak. Pelari yang hanya bersandar pada pengindraan karbon dioksida
akan berangkat dengan curah istirahat, menjalankan menit pertamanya
dengan \omterm{def:g10:cell-metabolism:fermentation}{fermentasi}, dan terpaksa melambat oleh asam laktat sebelum
pengantarannya menyusul.
\end{solution}

\begin{solution}{exo:g10:heart-lungs-effort:12}
$195 \times 0.110 \times 0.150 \approx \qty{3.2}{L/min}$ dan
$195 \times 0.160 \times 0.150 \approx \qty{4.7}{L/min}$. Volume
sekuncuplah yang diperbesar latihan: subjek yang kedua adalah yang
terlatih untuk ketahanan.
\end{solution}

\begin{solution}{exo:g10:heart-lungs-effort:13}
Oksigen yang dihirup: $0.21 \times 100 = \qty{21}{L/min}$; yang
dipakai: \qty{3.5}{L/min}, yaitu 17\%. Sisanya diembuskan lagi ---
udara embusan masih memuat sekitar 16\% oksigen, dan itulah sebabnya
pernapasan buatan dari mulut ke mulut berhasil.
\end{solution}

\begin{solution}{exo:g10:heart-lungs-effort:14}
Curahnya paling banyak sekitar $110 \times 0.12 \approx \qty{13}{L/min}$
alih-alih 24: $\dot V\!\mathrm{O_2}$maks-nya turun hampir separuh,
menjadi sekitar \qty{2}{L/min}. Pasiennya dapat berjalan dan bersepeda
perlahan tetapi tidak dapat mempertahankan usaha berat apa pun.
\end{solution}

\begin{solution}{exo:g10:heart-lungs-effort:15}
Pengantaran tiap liter turun seperempat, pengambilannya tetap secara
nisbi: $\dot V\!\mathrm{O_2}$maks turun 25\%. Sepanjang pekan-pekan
berikutnya tubuhnya membuat lebih banyak \omterm{def:g10:cells-common-unit:cell}{sel} darah merah, sehingga
oksigen yang terbawa tiap liter naik kembali menuju \qty{200}{mL}.
\end{solution}

\begin{solution}{pb:g10:heart-lungs-effort:1}
\textbf{1.} $62 \times 0.072 \approx 4.5$; $105 \times 0.100 = 10.5$;
$140 \times 0.116 \approx 16.2$; $175 \times 0.120 = 21.0$;
$198 \times 0.120 \approx \qty{23.8}{L/min}$.

\textbf{2.} $23.8/4.5 \approx 5.3$: frekuensinya menyumbang $198/62
\approx 3.2$, volume sekuncupnya $120/72 \approx 1.7$, dan $3.2 \times
1.7 \approx 5.3$.

\textbf{3.} Di atas sekitar \qty{225}{W} volume sekuncupnya tetap
\qty{120}{mL}; hanya denyut jantungnya yang menaikkan curahnya lebih jauh.

\textbf{4.} $220 - 17 = 203$: angka 198-nya berada dalam ketelitian aturannya.

\textbf{5.} $5/4.5 \approx \qty{1.1}{min}$, sekitar \qty{67}{s}; pada
\qty{300}{W}, $5/23.8 \approx \qty{0.21}{min}$, sekitar \qty{13}{s}.

\textbf{6.} $0.5 \times 13 = 6.5$; $1.2 \times 18 = 21.6$; $1.8 \times
24 = 43.2$; $2.3 \times 32 = 73.6$; $2.6 \times 44 \approx
\qty{114}{L/min}$.

\textbf{7.} $114/6.5 \approx 18$: volume tidalnya naik 5.2 kali,
frekuensinya 3.4 kali; volumenya menyumbang lebih banyak.

\textbf{8.} $2.6/5.0 = 52\%$ kapasitas vitalnya tiap embusan napas.

\textbf{9.} Tidak: \omterm{def:g10:heart-lungs-effort:ventilation}{ventilasinya} punya cadangan yang besar dan darah
yang meninggalkan paru-parunya tetap jenuh. Batasnya adalah volume
sekuncup jantungnya, yang menetapkan berapa banyak oksigen yang
diantarkan tiap denyut.

\textbf{10.} $0.7 \times 6.5 \approx \qty{4.6}{L/min}$ dan $0.7 \times
114 \approx \qty{80}{L/min}$.

\textbf{11.} $4.5 \times 0.052 \approx 0.23$; $10.5 \times 0.090
\approx 0.95$; $16.2 \times 0.120 \approx 1.94$; $21.0 \times 0.145
\approx 3.05$; $23.8 \times 0.155 \approx \qty{3.7}{L/min}$.

\textbf{12.} $3700/68 \approx \qty{54}{mL/min}$ tiap kilogram.
Frekuensinya sudah mencapai maksimumnya dan kenaikan dari 225 ke
\qty{300}{W} kecil: inilah $\dot V\!\mathrm{O_2}$maks-nya, milik
seorang murid yang bugar dan terlatih.

\textbf{13.} $3.7/0.23 \approx 16$: faktor jantungnya 5.3 dikalikan
faktor pengambilannya $155/52 \approx 3.0$.

\textbf{14.} $3.7 \times 20 = \qty{74}{kJ}$ tiap menit, sekitar
\qty{1230}{W}; dalam keadaan istirahat $0.23 \times 20/60 \approx \qty{77}{W}$;
di atas keadaan istirahat sekitar \qty{1150}{W}, jadi efisiensinya $300/1150 \approx
26\%$.

\textbf{15.} Kira-kira sebuah garis lurus dari \num{0.23} ke
\qty{3.05}{L/min} antara 0 dan \qty{225}{W}, lalu membengkok:
\qty{75}{W} terakhir hanya menambah \qty{0.65}{L/min}. Batas atasnya
sekitar \qty{3.7}{L/min}.

\textbf{16.} Antisipasi: harapan otaknya akan usaha itu diteruskan ke
pusat jantung dan kelenjar adrenalnya, sehingga frekuensinya naik
sebelum ada gerakan apa pun.

\textbf{17.} \omterm{def:g10:heart-lungs-effort:ventilation}{Alveolusnya} mengoksigenkan darahnya sepenuhnya bahkan pada
\omterm{def:g10:heart-lungs-effort:ventilation}{ventilasi} dan curah maksimal; paru-parunya bukan batasnya. Batasnya
adalah jumlah darah yang dapat dikirim jantungnya tiap menit.

\textbf{18.} $20/23.8 \approx 84\%$. Usus dan ginjalnya, yang tadinya
mengambil separuh curah istirahatnya, kini menerima sekitar
\qty{1}{L/min}; arterinya telah menyempit dan alirannya telah
dialihkan ke ototnya.

\textbf{19.} $198 \times 0.140 \times 0.155 \approx \qty{4.3}{L/min}$,
sekitar \qty{63}{mL/min} tiap kilogram.

\textbf{20.} \omterm{def:g10:heart-lungs-effort:output}{Curah jantungnya} berlipat 5.3, pengambilan oksigennya
berlipat 3, \omterm{def:g10:heart-lungs-effort:ventilation}{ventilasinya} berlipat 18; latihan mengubah volume sekuncup
jantungnya, dan bersamanya batas atasnya.
\end{solution}
